\begin{figure}[ht]
\centering
\begin{tikzpicture}[scale=1.0,
  node2/.style = {circle, draw=engblack, thick, fill=englime!25,
                  inner sep=0pt, minimum size=9mm, font=\small},
  ed/.style    = {thick, draw=enggray, ->, >=stealth},
  loop/.style  = {thick, draw=enggray, ->, >=stealth},
  elbl/.style  = {font=\scriptsize, fill=white, inner sep=1pt},
]
% Four nodes = 2-mers over {0,1}, placed on a square.
\node[node2] (00) at (0, 0)   {$00$};
\node[node2] (01) at (3, 0)   {$01$};
\node[node2] (10) at (0, 3)   {$10$};
\node[node2] (11) at (3, 3)   {$11$};

% Eight edges = 3-mers. Edge from ab to bc labelled abc.
\draw[loop] (00) to[out=200,in=140,looseness=8] node[elbl,left]{$000$} (00);
\draw[ed]   (00) -- node[elbl]{$001$} (01);
\draw[ed]   (01) -- node[elbl]{$011$} (11);
\draw[loop] (11) to[out=40,in=-20,looseness=8] node[elbl,right]{$111$} (11);
\draw[ed]   (11) -- node[elbl]{$110$} (10);
\draw[ed]   (10) -- node[elbl,pos=0.55]{$101$} (01);
\draw[ed]   (01) to[bend left=18] node[elbl,pos=0.4]{$010$} (10);
\draw[ed]   (10) -- node[elbl]{$100$} (00);
\end{tikzpicture}
\caption{The de Bruijn graph $B(2,3)$: each node is a binary
$2$-mer, and each directed edge is a $3$-mer joining the $2$-mer of
its first two bits to the $2$-mer of its last two. Every node has
in-degree and out-degree $2$, so by Euler's theorem the graph has an
Euler circuit that traverses all eight edges exactly once. Reading
the first bit of each edge along one such circuit
($000,001,011,111,110,101,010,100$) spells the cyclic de Bruijn
sequence $00011101$, which contains every one of the eight $3$-mers
exactly once. The same construction underlies shift-register
sequence design and the assembly of a genome from short overlapping
reads.}
\label{fig:vol02:ch17:de-bruijn}
\end{figure}
